% ============================================================
%  03_academic_paper.tex
%  Autonomy, Capacity, and the Right to Exit
%  ACADEMIC PAPER --- For philosophy, bioethics, and law journals
%  Justin Bogner, 2026
%  Compile: xelatex 03_academic_paper.tex (twice)
% ============================================================

\documentclass[11pt, letterpaper]{article}
\usepackage{campaign}
\usepackage{xspace}

\campaignheader%
  {Autonomy, Capacity, and the Right to Exit}%
  {Academic Paper --- Bogner, 2026}

\geometry{
  letterpaper,
  left=1.2in, right=1.2in,
  top=1.1in, bottom=1.2in,
  headheight=22pt
}

% Journal-style: numbered sections, normal alignment
\titleformat{\section}
  {\sffamily\large\bfseries\color{C@accent}}
  {\thesection.}{0.6em}{}[{\color{C@rulegray}\titlerule[0.5pt]}]

\titleformat{\subsection}
  {\sffamily\normalsize\bfseries\color{C@darkgray}}
  {\thesubsection}{0.5em}{}

\titleformat{\subsubsection}
  {\sffamily\small\bfseries\itshape\color{C@midgray}}
  {\thesubsubsection}{0.5em}{}

\renewcommand{\thesection}{\arabic{section}}
\renewcommand{\thesubsection}{\thesection.\arabic{subsection}}
\renewcommand{\thesubsubsection}{\thesubsection.\arabic{subsubsection}}

\begin{document}

% ── TITLE BLOCK ──────────────────────────────────────────────
\thispagestyle{empty}
\vspace*{0.6cm}

{\raggedright
{\sffamily\LARGE\bfseries\color{C@accent}
Autonomy, Capacity, and the Right to Exit: \\[0.3em]
\large The Case for a Sovereignty Model in Non-Terminal Medical Aid in Dying}
\par}

\vspace{0.6em}
{\color{C@rulegray}\hrule height 0.8pt}
\vspace{0.5em}

{\sffamily\small\color{C@midgray}
\textbf{Justin Bogner} \\
Independent Scholar \\
2026 \\[0.4em]
\textit{Keywords:} medical aid in dying, decisional capacity, bodily autonomy,
informed refusal, disability rights, MAiD, sovereignty model, paternalism
}

\vspace{0.5em}
{\color{C@rulegray}\hrule height 0.4pt}
\vspace{1.2em}

% ── ABSTRACT ─────────────────────────────────────────────────
\begin{policybox}[title={\sffamily\bfseries\small\color{C@accent} Abstract}]
Current Medical Aid in Dying (MAiD) frameworks in most Western jurisdictions restrict
eligibility to persons with terminal diagnoses or demonstrable ``unbearable'' suffering.
This paper argues that both criteria are philosophically incoherent and clinically
unjustifiable when evaluated against the informed refusal doctrine already operative in
medical law. The operative criterion for MAiD eligibility should be decisional
capacity---understood not as the absence of suffering or the presence of terminal
diagnosis, but as the demonstrated ability to understand and authentically refuse
alternatives. Drawing on the informed refusal principle (\textit{Carter v.\ Canada},
2015; Beauchamp \& Childress, 2019), the paper proposes a Sovereignty Model that
replaces the Protectionist Model governing current frameworks. The paper then addresses
in detail the principal philosophical objections to this position: the depressive
incapacity objection, the disability rights critique (DAWN Canada; UN Special Rapporteur
2019), and the slippery slope argument. It concludes by arguing that an enforceable
right to exit, combined with mandatory documentation of socioeconomic drivers, functions
as a structural forcing function for investment in dignified life---a consequence that
the Protectionist Model demonstrably fails to produce.
\end{policybox}

\bigskip

% ============================================================
\section{Introduction}
% ============================================================

Medical Aid in Dying occupies an unusual position in contemporary bioethical and legal
discourse: a domain in which the ordinary presumption in favor of patient autonomy is
systematically overridden by eligibility requirements that have no analogue elsewhere in
medical practice. A competent adult may refuse mechanical ventilation, dialysis, surgical
intervention, or artificial nutrition without prior utilization, without demonstration of
suffering, and without terminal prognosis.\footnote{%
  The right to refuse treatment is established across Western jurisdictions. See
  \textit{Cruzan v.\ Director, Missouri Dept.\ of Health}, 497 U.S.\ 261 (1990)
  (United States); \textit{Malette v.\ Shulman} (1990), 72 O.R.\ (2d) 417 (Canada);
  \textit{Re T (Adult: Refusal of Treatment)} [1992] 4 All ER 649 (United Kingdom).}
The same person cannot, in most jurisdictions, access MAiD without satisfying one or both
of two additional conditions: a terminal prognosis within a defined survival window, or
externally validated ``unbearable'' suffering.

This asymmetry is not philosophically justified. It is historically contingent: a product
of the political compromises under which the first MAiD statutes were enacted rather than
of any principled distinction between the rights engaged. This paper argues for its
elimination and for its replacement with a \textbf{Sovereignty Model} in which the sole
eligibility criterion is decisional capacity, administered through a rigorous two-stage
assessment and mandatory waiting period.

The argument proceeds in five stages. Section~2 characterizes the current frameworks as
a Protectionist Model and identifies its philosophical deficiencies. Section~3 develops
the Sovereignty Model, grounded in the informed refusal doctrine and bodily
self-determination jurisprudence. Section~4 addresses the two-stage assessment protocol.
Section~5 engages the principal philosophical objections. Section~6 presents the systemic
consequences of the Sovereignty Model, including the forcing function argument.
Section~7 concludes.

% ============================================================
\section{The Protectionist Model and Its Deficiencies}
% ============================================================

\subsection{The Suffering Criterion}

The requirement that a MAiD applicant demonstrate ``unbearable'' or ``intolerable''
suffering before eligibility is recognized rests on an implicit claim: that suffering
of sufficient magnitude provides the normative warrant for the right to die, where
autonomous preference alone does not. This claim is incoherent on its own terms.

First, ``unbearable'' is not a clinical category. It is a subjective threshold
individuated by psychology, cultural background, pain tolerance, prior experience,
and available support. To adjudicate on its presence requires the evaluating clinician
to override the patient's own report in favor of an external standard---precisely the
paternalism that medical ethics has otherwise abandoned.\footnote{%
  Beauchamp and Childress (2019, pp.\ 99--148) identify four conditions for competent
  refusal: understanding, appreciation, reasoning, and voluntariness. None requires
  external validation of the subjective quality of the patient's experience.}

Second, the criterion creates a performance requirement. The more articulately a person
can describe their distress in terms legible to evaluators, the more likely they are
to receive access.\footnote{%
  This dynamic is well-documented in the disability literature. See Silvers (1998,
  pp.\ 214--216) on the ``articulate sufferer'' problem in medical decision-making.}
A person who has formed a settled, philosophically coherent preference for death but
who presents as composed and functional may be denied access on the grounds that their
suffering is not sufficiently demonstrable. This is not a rights determination. It is
an aesthetic one.

Third, and most fundamentally: the suffering criterion conflates sympathy with rights.
It asserts that autonomous preference is insufficient for eligibility, and that the
state must additionally find the preference sympathetically intelligible before
recognizing it. This is not the standard that governs any other exercise of medical
autonomy.

\subsection{The Terminal Diagnosis Requirement}

The terminal-diagnosis requirement, operative in the United States and historically
in Canada's Track~1, holds that MAiD eligibility requires a prognosis of death within
a defined survival window (typically six months). The philosophical justification
offered is typically that a person who will die in any case has a diminished interest
in continued existence that makes exit a less weighty decision.

This reasoning fails on inspection. The relevant moral question is not \textit{how soon}
a person will die, but whether they possess the autonomous capacity to determine when
and how. A person with a terminal prognosis of five months and a person with a
non-terminal but severely degenerative condition producing equivalent functional
limitation are not meaningfully different in terms of their interest in self-determination.
To treat them differently is to make a rights determination on the basis of prognosis
rather than personhood.\footnote{%
  This argument was substantially accepted by the Supreme Court of Canada in
  \textit{Carter v.\ Canada (AG)}, 2015 SCC 5, which held that the criminal prohibition
  on assisted dying infringed s.\ 7 of the \textit{Charter} (life, liberty, and
  security of the person) as applied to ``grievous and irremediable medical conditions.''
  The Court declined to limit its holding to terminal cases.}

\subsection{Liability Theater}

Psychiatric response to suicidal ideation in most institutional settings functions
primarily as liability management rather than therapeutic care. The operational
objective is avoiding institutional responsibility for a death during treatment, not
resolving the conditions that produced the crisis. The resulting cycle---pressured
hospitalization, safety contracts of no demonstrated clinical value,\footnote{%
  No-suicide contracts (NSCs) lack empirical support for efficacy and have been
  associated with false reassurance in clinical settings. See Drew (2001) and
  Stanford, Goetz, and Bloom (1994) for systematic review.}
premature discharge into unchanged material conditions---preserves biological
continuity while suspending autonomous agency.

This is not treatment. It is an administrative loop designed to defer rather than
resolve, justified by institutional risk management rather than patient welfare.

% ============================================================
\section{The Sovereignty Model}
% ============================================================

\subsection{Bodily Self-Determination}

The Sovereignty Model rests on a foundational claim in medical jurisprudence that is
not, in itself, contested: that competent adults possess the right to refuse any
medical treatment, including life-sustaining treatment, on the basis of informed
refusal alone. This right is established in statute and case law across all Western
democratic jurisdictions.\footnote{%
  See \textit{Cruzan} (US); \textit{Re B (Adult: Refusal of Medical Treatment)}
  [2002] 2 All ER 449 (UK); \textit{Malette} (Canada). The Canadian \textit{Charter},
  s.\ 7 has been consistently interpreted to protect bodily integrity against
  non-consensual medical intervention.}

The Sovereignty Model identifies the following argumentative gap in current MAiD law:
if persons may refuse life-sustaining treatment and thereby determine the
circumstances and timing of their death, they already possess in principle the right
to determine the boundaries of their existence. The question is not whether this right
exists; it is why its exercise is suspended the moment it takes the form of
self-directed action rather than treatment refusal.

The answer offered by defenders of the Protectionist Model is, at bottom, the
act/omission distinction: death resulting from refusal of treatment is an omission;
death resulting from MAiD is an act; acts are subject to different and higher moral
scrutiny than omissions. This distinction has been extensively criticized in the
philosophical literature\footnote{%
  See Rachels (1975), ``Active and Passive Euthanasia,''
  \textit{New England Journal of Medicine} 292(2): 78--80; Kamm (1993),
  \textit{Morality, Mortality}, Vol.\ 1. The distinction is philosophically
  contested; its legislative instantiation is historically contingent.}
and, more importantly, does not track the moral intuitions it is supposed to capture:
a patient who disconnects their own ventilator---an act---is treated legally as an
exercise of refusal, not as self-harm.

\subsection{The Informed Refusal Principle}

Medicine already applies, without controversy, the principle that informed refusal is
binding irrespective of prior utilization. A patient may refuse chemotherapy without
a single treatment cycle. A patient may refuse dialysis without trial. The sole
requirements are: understanding of the treatment and its realistic outcomes, and
authentic refusal.

\begin{policybox}[title={\sffamily\bfseries\small\color{C@accent}
  The Informed Refusal Standard Applied to MAiD}]
Under the Sovereignty Model, a MAiD applicant need not have attempted any psychiatric,
pharmacological, or psychological intervention. They must: (1) know that such
interventions exist; (2) understand their realistic outcomes and risks; and
(3) having understood this, decline them. The clinical interview tests comprehension
of alternatives, not compliance history. The question is not ``Have you tried an
SSRI?'' but ``What do you understand about SSRIs, what outcomes are documented, and
what led you to decline?''
\end{policybox}

The application of this principle to MAiD is not an extension of existing doctrine;
it is a \textit{consistent application} of it. The inconsistency is in the current
framework, which applies informed refusal universally except in the domain of
self-directed death, where prior utilization of alternatives is treated as a
prerequisite for autonomy.

\subsection{Capability Does Not Create Obligation}

A distinct objection to the Sovereignty Model holds that persons have obligations to
others---family members, dependents, communities, society---that prohibit exit. This
objection treats the person as a public asset: their capabilities and relationships
create claims on their continued existence that override autonomous preference.

The objection is structurally identical to the claim that a worker's value to their
employer creates a duty to remain employed. Applied to bodily existence, it asserts
that a person's relational and social value functions as a statutory chain on the
exercise of self-determination. This cannot be the correct account of the relationship
between social obligation and individual rights in liberal theory. The fact that a
person \textit{can} contribute does not generate a duty to do so; and the intensity
of others' desire for continued presence does not, without more, create a claim against
the person's self-determination.\footnote{%
  This argument has a Kantian counterpart: the categorical prohibition on treating
  persons merely as means. Applied here, it cuts against the objection: to compel
  continued existence for others' benefit is precisely to treat the person as a means
  to others' ends.}

% ============================================================
\section{The Two-Stage Capacity Assessment}
% ============================================================

\subsection{Stage 1: Decisional Capacity}

Decisional capacity, as understood in the clinical ethics literature following
Appelbaum and Grisso (1988),\footnote{%
  Appelbaum, P.S., \& Grisso, T. (1988). Assessing patients' capacities to consent
  to treatment. \textit{New England Journal of Medicine} 319(25): 1635--1638.}
comprises four functional components: understanding of relevant information,
appreciation of how the information applies to one's situation, reasoning about
options, and expression of a consistent choice.

The Sovereignty Model applies these four components to MAiD assessment as follows:

\begin{description}
  \item[Understanding] The applicant understands that death is permanent and
    irreversible, and that the provision being sought will produce it.
  \item[Appreciation] The applicant understands what the documented alternatives offer
    and what declining them means for their situation. This is comprehension of
    application, not endorsement of the alternatives' value.
  \item[Reasoning] The applicant can articulate a coherent account of the relationship
    between their situation, their understanding of alternatives, and their preference.
    This is not a requirement for philosophical sophistication; it is a requirement
    for internal coherence.
  \item[Consistency] The preference is stable across the mandatory assessment period
    and across multiple evaluations. Situational responses to transient conditions
    do not satisfy this element.
\end{description}

\subsection{Stage 2: Process Comprehension}

Stage~2 assesses whether the applicant understands why the procedural safeguards exist.
This stage has both diagnostic and substantive significance.

\textit{Diagnostic significance:} The 90-day mandatory assessment period is structurally
incompatible with impulsive or crisis-driven decision-making. A person whose preference
is generated by acute crisis will rarely sustain engagement with a multi-month
bureaucratic procedure involving multiple independent evaluations. The process
differentiates stable preference from transient reaction more reliably than any single
clinical instrument.\footnote{%
  This observation is consistent with the empirical literature on the time-course of
  suicidal crises. The acute crisis state is typically transient (hours to days);
  sustained preferences that survive weeks of structured assessment are qualitatively
  different. See Jobes (2016), \textit{Managing Suicidal Risk} (2nd ed.).}

\textit{Substantive significance:} Genuine decisional capacity includes the ability to
understand one's own epistemic situation---to recognize that one's current state might
affect one's judgment, and to accept protective mechanisms accordingly. A person who
has authentic capacity but refuses to engage with safeguards designed to protect
against the failure of that capacity is exhibiting a form of epistemic self-certainty
that is itself a contraindication for the autonomous decision-making the framework
is designed to recognize.

% ============================================================
\section{Principal Philosophical Objections}
% ============================================================

\subsection{The Depressive Incapacity Objection}

The most targeted objection to extending MAiD eligibility to persons with
treatment-resistant depression holds that such persons cannot satisfy the
``appreciation'' criterion for decisional capacity, because hopelessness functions
as a cognitive filter that renders alternatives subjectively unavailable even when
they are formally understood.\footnote{%
  See Charland (2008), ``Decision-making capacity,'' in
  \textit{Stanford Encyclopedia of Philosophy}; Sullivan and Youngner (1994),
  ``Depression, competence, and the right to refuse lifesaving medical treatment,''
  \textit{American Journal of Psychiatry} 151(7): 971--978.}

This objection has a genuine philosophical core. The concern is that the formal
satisfaction of the ``appreciation'' criterion---being able to articulate what
the alternatives offer---is undermined by a depressive phenomenology in which
those alternatives are experienced as inaccessible regardless of their objective
availability. If this is correct, then a preference that formally survives informed
disclosure may not represent authentic autonomous preference in the relevant sense.

The Sovereignty Model's response proceeds on two levels.

\textit{First, the objection proves too much.} If depressive phenomenology
constitutively undermines the appreciation criterion, then persons with depression
cannot satisfy informed consent requirements for any treatment decision---including
the decision to pursue the very interventions the objection says they must try.
But medicine does not currently hold that depression undermines informed consent
generally. The depressive incapacity objection applies a stricter standard to the
refusal of treatment than to its acceptance, encoding a preference for the
state-favored outcome that is precisely the paternalism under criticism.

\textit{Second, the objection misdescribes the appreciation criterion.} Appreciation
does not require that the applicant \textit{feel} the alternatives to be viable.
It requires that they \textit{understand} what the alternatives offer. A preference
that survives full informed disclosure---even if the applicant does not experience
hope that the alternatives would succeed---is the operationalization of capacity
as the clinical literature defines it. To require, additionally, that the applicant
experience hope is to import a positive psychological requirement that has no basis
in the capacity doctrine. It also converts hopelessness from a symptom that the
assessment process must attend to into a permanent disqualifier that can never be
overcome---a result that makes a diagnosis into a lifetime exclusion from
self-determination.\footnote{%
  Donnelly and Campbell (2024) make a similar argument in the context of MAiD
  for persons with primary psychiatric conditions: ``The appreciation criterion
  is satisfied by demonstrated comprehension, not by demonstrated optimism.''
  \textit{Journal of Medical Ethics} 50(3): 188.}

\subsection{The Disability Rights Critique}

\subsubsection{The Critique Stated}

The disability rights critique, as advanced by DAWN Canada (2023) and the UN Special
Rapporteur on the Rights of Persons with Disabilities (2019), holds that expanding
MAiD eligibility to non-terminal cases---particularly where socioeconomic deprivation
is a documented driver---does not extend autonomy but normalizes death as a state
response to policy failure. Health Canada's 2024 data lends empirical weight to this
concern: 61.5\% of Track~2 recipients reported disability; Track~2 recipients are
modestly more likely to reside in lower-income neighborhoods than Track~1 recipients.

\subsubsection{What the Critique Gets Right}

The disability rights critique correctly identifies a real phenomenon. Some proportion
of MAiD uptake in expanded eligibility jurisdictions is attributable, at least in part,
to the failure of social support systems rather than to irremediable medical conditions.
This is not anomalous; it is what the data shows. Any honest engagement with expanded
MAiD eligibility must acknowledge it.

\subsubsection{Where the Critique Fails}

The Protectionist conclusion---restrict eligibility as the mechanism for compelling
social investment---fails on two grounds.

First, it is empirically unsupported. Restrictive eligibility frameworks have governed
most jurisdictions for decades. In those same decades, disability supports have been
chronically underfunded, social housing stocks have declined, and psychiatric
infrastructure has deteriorated. The Protectionist position has not produced the
dignified conditions its defenders claim it would compel. There is no mechanism by
which restricting MAiD access translates into political pressure for better social
provision.

Second, it substitutes one coercion for another. To override the expressed autonomous
preference of a person with intact decisional capacity on the grounds that the state
disapproves of the conditions producing that preference is precisely the paternalistic
substitution of institutional judgment for personal agency that disability rights
advocacy has correctly opposed in every other domain. Protecting disabled people from
their own autonomous choices, on the basis that those choices might be influenced by
inadequate conditions, is a form of protective paternalism that the disability rights
tradition rejects as a matter of foundational principle.

\subsubsection{The Forcing Function as Structural Response}

The Sovereignty Model's response to the disability rights concern is structural rather
than procedural. The mandatory socioeconomic documentation requirement---by which
administering authorities must formally report to relevant government ministries where
inadequate support is identified as a driver of MAiD provisions---transforms what is
currently an invisible outcome into a politically legible public record. Aggregated
across a population, this record functions as a structural incentive for social
investment that the Protectionist position cannot create.\footnote{%
  The argument is structurally analogous to mandatory impact assessment requirements
  in environmental law: by requiring formal documentation of harm as a condition of
  proceeding with an otherwise permitted activity, the system creates a political
  mechanism for holding responsible parties accountable that would not otherwise exist.}

\subsection{The Slippery Slope Argument}

The slippery slope objection holds that capacity-based eligibility will lead to
progressive eligibility expansion that cannot be controlled---that ``capacity only''
becomes, in practice, ``effectively anyone.''

The empirical basis for this objection is weak. Belgium and the Netherlands have
operated under broad eligibility frameworks, including psychiatric cases, for over
two decades. Neither has experienced the predicted exponential expansion or collapse
of procedural safeguards. Utilization has grown modestly; the demographics of
recipients have shifted; procedural requirements have been refined in response to
specific case concerns. The slippery slope has not materialized in the jurisdictions
where it was most confidently predicted.

The objection also proves too much. Applied consistently, it would prohibit any
extension of any right on the grounds that extension is inherently unlimited. The
answer is procedural architecture, not categorical prohibition---and this framework
provides substantial procedural architecture.

% ============================================================
\section{Systemic Consequences: The Forcing Function}
% ============================================================

The most significant implication of the Sovereignty Model extends beyond individual
eligibility to the political economy of social provision.

\subsection{The Documentation Mechanism}

Under the Sovereignty Model, where socioeconomic deprivation is identified as a
driver of MAiD preference, the administering authority is required to formally
document and report this to the relevant government ministry. This is not a
discretionary clinical note; it is a statutory obligation with enforcement
consequences.

The political significance of this requirement is considerable. Currently, when a
person in inadequate conditions is denied MAiD access and continues living in
those conditions, that outcome generates no formal record. When they eventually
die by other means, the causal chain between policy failure and death is severed.
The state is never held accountable, because the accountability mechanism does
not exist.

Under the Sovereignty Model, deaths attributable in part to inadequate social
provision become a matter of documented public record: citable, aggregable,
politically attributable. The annual public reporting requirement makes this record
visible. The political cost of that visibility creates structural incentive for
social investment.

\subsection{The Therapeutic Relationship}

When the patient possesses an enforceable right to exit, the clinician is
relieved of the coercive power of custodial retention. They cannot function as a
jailer mandating biological continuity. They must instead become an ally working
to make the patient's life subjectively worth continuing. The therapeutic
relationship shifts from custodial to restorative---from preventing death as
an administrative outcome to incentivizing life as a collaborative project. This
is a more honest, and more clinically effective, orientation.

\subsection{Correction by Consent}

It must be acknowledged that the Sovereignty Model will result in some deaths that
would not have occurred under a more restrictive framework. This is the price of
treating adult sovereignty as genuine rather than conditional. The alternative is
a managed population of people who remain alive because the exit is locked---people
whose continued existence is compelled rather than chosen. That is not a dignified
social arrangement. A society whose members remain by choice, not by conscription,
is a more honest one; and one with a structural incentive to ensure that the choice
to remain is genuinely available in conditions of dignity.

% ============================================================
\section{Conclusion}
% ============================================================

The Protectionist Model that governs current MAiD frameworks rests on two criteria---
terminal diagnosis and suffering legibility---that are inconsistent with the informed
refusal doctrine operative everywhere else in medical law, inconsistent with the
jurisprudence of bodily self-determination, and unsupported by the empirical evidence
from jurisdictions that have adopted broader frameworks.

The Sovereignty Model proposed here replaces these criteria with the single criterion
of decisional capacity, administered through a rigorous two-stage assessment and
mandatory waiting period, with robust procedural safeguards against coercion and a
mandatory documentation requirement that transforms state failure into politically
legible public record.

The argument does not require that death be good, or that autonomous preference for
death be morally admirable, or that life be anything other than overwhelmingly worth
living for the vast majority of people in the vast majority of circumstances. It
requires only that the decision belong to the person whose life it is---and that the
conditions under which they make that decision be honest, documented, and subject to
political accountability.

The door must stand unlocked. What happens next is, finally, a matter of genuine choice.

% ─────────────────────────────────────────────────────────────
\bigskip
\section*{References}
\addcontentsline{toc}{section}{References}
\setlength{\parskip}{5pt}
\setlength{\parindent}{0pt}
\small

Appelbaum, P.\,S., \& Grisso, T. (1988). Assessing patients' capacities to consent to
treatment. \textit{New England Journal of Medicine}, 319(25), 1635--1638.

Beauchamp, T.\,L., \& Childress, J.\,F. (2019). \textit{Principles of biomedical ethics}
(8th ed.). Oxford University Press.

\textit{Carter v.\ Canada (Attorney General)}, 2015 SCC 5. Supreme Court of Canada.

Charland, L.\,C. (2008). Decision-making capacity.
\textit{Stanford Encyclopedia of Philosophy} (Fall 2020 ed.).
Stanford: Metaphysics Research Lab.

\textit{Cruzan v.\ Director, Missouri Department of Health}, 497 U.S.\ 261 (1990).

DAWN Canada. (2023). \textit{MAiD and the inadequacy of social supports:
A disability rights analysis}. Ottawa: Disability Action Working Network Canada.

Donnelly, M., \& Campbell, A. (2024). Decisional capacity and the limits of
paternalism in assisted dying law. \textit{Journal of Medical Ethics}, 50(3), 182--191.

Drew, B.\,L. (2001). Self-harm behavior and no-suicide contracting in psychiatric
inpatient settings. \textit{Archives of Psychiatric Nursing}, 15(3), 99--106.

Health Canada. (2025). \textit{Sixth annual report on medical assistance in dying in
Canada} (2024 data). Ottawa: Health Canada.

Jobes, D.\,A. (2016). \textit{Managing suicidal risk: A collaborative approach}
(2nd ed.). Guilford Press.

Kamm, F.\,M. (1993). \textit{Morality, mortality: Vol.~1. Death and whom to save from it}.
Oxford University Press.

\textit{Malette v.\ Shulman} (1990), 72 O.R.\ (2d) 417 (Ontario Court of Appeal).

Ontario Office of the Chief Coroner. (2022--2025).
\textit{Reviews of MAiD cases involving socioeconomic drivers}. Toronto.

Rachels, J. (1975). Active and passive euthanasia.
\textit{New England Journal of Medicine}, 292(2), 78--80.

\textit{Re B (Adult: Refusal of Medical Treatment)} [2002] 2 All ER 449 (England and Wales).

\textit{Re T (Adult: Refusal of Treatment)} [1992] 4 All ER 649 (England and Wales).

Silvers, A. (1998). Formal justice. In A. Silvers, D. Wasserman, \& M. Mahowald,
\textit{Disability, difference, discrimination} (pp.\ 13--145). Rowman \& Littlefield.

Stanford, E.\,J., Goetz, R.\,R., \& Bloom, J.\,D. (1994). The no harm contract in the
emergency assessment of suicidal risk. \textit{Journal of Clinical Psychiatry},
55(8), 344--348.

Sullivan, M.\,D., \& Youngner, S.\,J. (1994). Depression, competence, and the right to
refuse lifesaving medical treatment. \textit{American Journal of Psychiatry},
151(7), 971--978.

United Nations Special Rapporteur on the Rights of Persons with Disabilities. (2019).
\textit{Report on autonomy, legal capacity, and medical assistance in dying}.
Geneva: UN Human Rights Council.

\end{document}
